\section{Conclusion}
This study evaluates budgeted reobservation at two levels: paired building predictions and online disaster-inspection answers. In the recorded imagery setting, a smaller camera footprint changes the perception endpoint and enables useful selection under a fixed building budget. Calibrated entropy improves macro-F1 over the recorded random allocation, but does not outperform raw entropy, and its improvement does not extend uniformly to probability-quality metrics. The online experiment shows a different result: additional observations produce little change in question correctness and are often concentrated on already-correct damage questions.

The supported conclusion is that observation change, uncertainty change, and task value must be measured separately. A functioning feedback loop and higher answer confidence are insufficient evidence of successful sensing allocation. The present results remain conditional on the exposed backbone and archived execution protocol. Future evaluation should first restore full provenance and budget matching, then test whether task-conditioned evidence requirements improve the allocation of new observations on genuinely unseen events. No such improvement is claimed by the current experiments.
